% ──────────────────────────────────────────────────────────────────────────────
%  Section 7 — Conclusion
% ──────────────────────────────────────────────────────────────────────────────

\section{Conclusion}
\label{sec:conclusions}

The closed-form expression for the peak-to-peak modular spread
$\Delta_{\max}(S,\ell)$ delivers a single operational reading: the
actual electrical state of a terminal cannot be determined from the
phase current alone, whether read as a single-phase modulus or as the
full three-phase set. The reason is structural. The line is
geometrically not circulant, and under balanced positive-sequence
voltage at both terminals the natural asymmetry of the per-unit-length
matrices has no degree of freedom left in the voltage state; it is
forced into the current and into the per-phase split of $P$ and $Q$.
The current is therefore unbalanced by construction, set by the tower
geometry rather than by a fault or an asymmetric excitation. At the
anchor of \autoref{sec:example} this materialises as a
$\SI{9.08}{\percent}$ peak-to-peak spread and a
$\SI{6.13}{\percent}$ negative-sequence ratio, both produced entirely
by geometry.

In any actual deployment the voltages at the two ends are not
perfectly balanced either, so the asymmetry redistributes across all
four signals -- voltage, current, active and reactive power per phase
-- in proportions that depend on the loading direction and on the
stiffness of each bus. A current-only reading, even on all three
phases, cannot recover that distribution. The actual state of the
terminal is the four-signal package $(V,I,P,Q)$ per phase, with phase
angles, measured at both ends.
